% P11 paper module: R21 exact cross-polar defect and relative gauge reparametrization.

\subsection{Cross-polar defect and exact relative gauge reparametrization}
\label{subsec:o3t-crosspolar}

Keep the odd-sector notation at fixed $0<R<S<T_0<U$ from
Subsections~\ref{subsec:o3n-gauge-intertwining},
\ref{subsec:o3r-sqrt-offblock}, and~\ref{subsec:o3s-relative-polar}.  Thus
\[
 W:=W_{R,S,-}^{[T_0]},
 \qquad
 A_R:=A_{T_0,U}^{R,-},
 \qquad
 A_S:=A_{T_0,U}^{S,-},
\]
\[
 X_R=U_RA_R^{1/2},
 \qquad
 X_S=U_SA_S^{1/2},
\]
and the genuine relative polar defect is
\[
 \mathfrak P_U:=U_SW-WU_R.
\]
The purpose of this subsection is to express this defect through a single cross-product
of the two metric polar products.  No separate estimate of the individual commutators
$[G_{R,T_0},G_{R,U}]$ and $[G_{S,T_0},G_{S,U}]$ is used.

\begin{proposition}[Exact cross-polar identity]
\label{prop:o3t-crosspolar}
Define
\[
 \boxed{
 Z_U
 :=A_S^{-1/2}X_S^*WX_RA_R^{-1/2}.
 }
\]
Then
\[
 \boxed{
 Z_U=U_S^*WU_R.
 }
\tag{O3T.1}\label{eq:o3t-Z}
\]
In particular $Z_U$ is an isometry from the $R$-source space into the $S$-source
space.
\end{proposition}

\begin{proof}
Since $X_S^*=A_S^{1/2}U_S^*$ and $X_R=U_RA_R^{1/2}$,
\[
\begin{aligned}
 Z_U
 &=A_S^{-1/2}(A_S^{1/2}U_S^*)W(U_RA_R^{1/2})A_R^{-1/2}\\
 &=U_S^*WU_R.
\end{aligned}
\]
Both polar factors are unitary and $W$ is an isometry, hence so is $Z_U$.
\end{proof}

\begin{proposition}[Exact isometric gauge-defect reparametrization]
\label{prop:o3t-gauge-defect}
One has the exact operator identity
\[
 \boxed{
 Z_U-W=-U_S^*\mathfrak P_U.
 }
\tag{O3T.2}\label{eq:o3t-gauge-defect}
\]
Consequently, for every source vector $f$,
\[
 \boxed{
 \|(Z_U-W)f\|=\|\mathfrak P_Uf\|,
 }
\tag{O3T.3}\label{eq:o3t-vector-norm}
\]
and therefore
\[
 \boxed{
 Z_U-W\xrightarrow[s]{}0
 \quad\Longleftrightarrow\quad
 \mathfrak P_U\xrightarrow[s]{}0.
 }
\tag{O3T.4}\label{eq:o3t-strong-equivalence}
\]
In operator norm,
\[
 \boxed{
 \|Z_U-W\|=\|\mathfrak P_U\|,
 }
\tag{O3T.5}\label{eq:o3t-norm-equality}
\]
so norm convergence is equivalent as well.
\end{proposition}

\begin{proof}
Equation~\eqref{eq:o3t-Z} gives
\[
 Z_U-W
 =U_S^*WU_R-U_S^*U_SW
 =-U_S^*(U_SW-WU_R)
 =-U_S^*\mathfrak P_U.
\]
Unitary invariance of the vector and operator norms yields
\eqref{eq:o3t-vector-norm} and~\eqref{eq:o3t-norm-equality}.  The pointwise norm
identity gives the strong-convergence equivalence~\eqref{eq:o3t-strong-equivalence}.
\end{proof}

The preceding proposition is an exact reparametrization rather than an estimate: no
information about the relative polar defect is lost in passing from $\mathfrak P_U$ to
$Z_U-W$.  The same information can be written before the inverse-square-root
normalization as follows.

\begin{corollary}[Unnormalised relative cross-polar defect]
\label{cor:o3t-Crel}
Put
\[
 \boxed{
 \mathcal C_U^{\rm rel}
 :=X_S^*WX_R-A_S^{1/2}WA_R^{1/2}.
 }
\tag{O3T.6}\label{eq:o3t-Crel}
\]
Then
\[
 \boxed{
 \mathcal C_U^{\rm rel}
 =A_S^{1/2}(Z_U-W)A_R^{1/2}.
 }
\tag{O3T.7}\label{eq:o3t-Crel-factor}
\]
Hence
\[
 \boxed{
 \|\mathfrak P_U\|
 \ge
 \frac{\|\mathcal C_U^{\rm rel}\|}
 {\sqrt{\|A_S\|}\sqrt{\|A_R\|}},
 }
\tag{O3T.8}\label{eq:o3t-Crel-lower}
\]
while
\[
 \boxed{
 \|\mathfrak P_U\|
 \le
 \|A_S^{-1/2}\|\,
 \|\mathcal C_U^{\rm rel}\|\,
 \|A_R^{-1/2}\|.
 }
\tag{O3T.9}\label{eq:o3t-Crel-upper}
\]
For the concrete P11 family, Lemma~\ref{lem:o3l-relative-upper} gives
$\|A_R\|,\|A_S\|\ll e^U/U$, and therefore
\[
 \boxed{
 \|\mathfrak P_U\|
 \gtrsim
 \frac{U}{e^U}\,\|\mathcal C_U^{\rm rel}\|.
 }
\tag{O3T.10}\label{eq:o3t-P11-weighted-lower}
\]
The constants here may depend on the fixed levels $R,S,T_0$.
\end{corollary}

\begin{proof}
Using~\eqref{eq:o3t-Z},
\[
 X_S^*WX_R
 =A_S^{1/2}Z_UA_R^{1/2},
\]
which proves~\eqref{eq:o3t-Crel-factor}.  The first norm inequality follows from
\[
 \|\mathcal C_U^{\rm rel}\|
 \le
 \sqrt{\|A_S\|}\,\|Z_U-W\|\,\sqrt{\|A_R\|}
\]
and~\eqref{eq:o3t-norm-equality}.  Conversely,
\[
 Z_U-W
 =A_S^{-1/2}\mathcal C_U^{\rm rel}A_R^{-1/2},
\]
which gives~\eqref{eq:o3t-Crel-upper}.  Finally insert the relative-metric upper bounds
from Lemma~\ref{lem:o3l-relative-upper} into~\eqref{eq:o3t-Crel-lower}.
\end{proof}

There is also a source-space compression which is unitarily conjugate to the R15 gauge
compression.

\begin{corollary}[Cross-polar compression of the gauge]
\label{cor:o3t-Omega}
Define
\[
 \boxed{
 \Omega_U:=Z_U^*W
 =A_R^{-1/2}X_R^*W^*X_SA_S^{-1/2}W.
 }
\]
Then
\[
 \boxed{
 \Omega_U=U_R^*\Gamma_UU_R,
 }
\tag{O3T.11}\label{eq:o3t-Omega}
\]
where $\Gamma_U=W^*U_SWU_R^*$.  Consequently
\[
 \boxed{
 \|\Omega_U-I\|=\|\Gamma_U-I\|.
 }
\tag{O3T.12}\label{eq:o3t-Omega-norm}
\]
\end{corollary}

\begin{proof}
By~\eqref{eq:o3t-Z},
\[
 \Omega_U
 =U_R^*W^*U_SW
 =U_R^*(W^*U_SWU_R^*)U_R
 =U_R^*\Gamma_UU_R.
\]
The operator-norm identity follows by unitary invariance.
\end{proof}

\begin{remark}[No automatic strong-convergence transfer through $\Omega_U$]
\label{rem:o3t-Omega-strong-firewall}
Equation~\eqref{eq:o3t-Omega} is an exact unitary conjugacy at each fixed $U$, and
\eqref{eq:o3t-Omega-norm} is therefore exact.  Since the conjugating unitary $U_R$
depends on $U$, one must not infer from this identity alone an equivalence between
strong convergence of $\Omega_U$ and strong convergence of $\Gamma_U$.  The strong
statement that is exact without an additional hypothesis is instead
\eqref{eq:o3t-strong-equivalence} for $Z_U-W$ and $\mathfrak P_U$.
\end{remark}

\begin{remark}[R21-D scaling firewall]
\label{rem:o3t-scale-firewall}
The raw size of $\mathcal C_U^{\rm rel}$ is not itself a scale-invariant gauge measure.
Indeed, at fixed base metrics, simultaneously rescale the two future metrics by the
same scalar $t>0$.  Then
\[
 A_R,A_S\mapsto tA_R,tA_S,
 \qquad
 X_R,X_S\mapsto\sqrt t\,X_R,\sqrt t\,X_S,
\]
while the polar factors $U_R,U_S$, hence $Z_U-W$ and $\mathfrak P_U$, are unchanged.
In contrast,
\[
 \mathcal C_U^{\rm rel}\mapsto t\mathcal C_U^{\rm rel}.
\]
Thus a large unnormalised cross-polar norm cannot be promoted to a large gauge defect
without the metric normalization in~\eqref{eq:o3t-Crel-lower}--\eqref{eq:o3t-Crel-upper}.
\end{remark}

\begin{remark}[R21 countermodel firewall]
\label{rem:o3t-countermodels}
The cross-polar defect passes the canonical inclusion countermodels that blocked the
previous routes:
\begin{center}
\begin{tabular}{@{}p{0.23\textwidth}p{0.46\textwidth}p{0.21\textwidth}@{}}
\toprule
Model & Existing phenomenon & Cross-polar verdict \\
\midrule
R14, Proposition~\ref{prop:o3m-countermodel}
& $Q=W$ and $\Theta=0$, but the actual polar gauge is nontrivial
& $Z_U-W\ne0$ \\
R19, Proposition~\ref{prop:o3r-polar-countermodel}
& nonzero/arbitrarily large modulus and square-root leakage, but $U_S=I$, $U_R=I$
& $Z_U-W=0$ \\
R20 matched, Proposition~\ref{prop:o3s-matched-countermodel}
& nontrivial individual polar factors and arbitrarily large raw commutators, but $U_SW=WU_R$
& $Z_U-W=0$ \\
R20 combined, Proposition~\ref{prop:o3s-combined-countermodel}
& modulus leakage plus large individual commutators, still $U_SW=WU_R$
& $Z_U-W=0$ \\
\bottomrule
\end{tabular}
\end{center}
These conclusions are not separate model calculations: they follow immediately from
the exact identity~\eqref{eq:o3t-gauge-defect}.  Thus $Z_U-W$ rejects precisely the
matched polar configurations that should vanish and detects the R14 gauge-active
configuration.
\end{remark}

\begin{openproblem}[R21-F: concrete P11 cross-polar asymptotics]
\label{open:o3t-crosspolar}
Determine the asymptotic behavior of
\[
 Z_U-W
 =A_S^{-1/2}X_S^*WX_RA_R^{-1/2}-W
\]
for the concrete P11 metric family.  For the strong terminal problem the relevant task
is fixed-vector control of $(Z_U-W)f$; an operator-norm lower bound alone would not rule
out strong convergence.  Conversely, a norm upper bound tending to zero would imply
strong gauge compatibility.  Equivalently one may estimate
$\mathcal C_U^{\rm rel}$ together with sufficiently sharp control of the positive and
inverse square roots of $A_R,A_S$.

The present fixed-pair rank-one asymptotics do not provide this control:
Proposition~\ref{prop:o3n-rankone-insufficiency} shows that the inverse-square-root scale
is not determined by the existing leading Gram data, and Remark~\ref{rem:o3p-firewall}
identifies the remaining problem as operator-valued orientation of positive/inverse
square roots under $R\to S$.  Hence no conclusion about
\[
 \Gamma_U\to I,
 \qquad
 K_{R,S}^{T_0,U}\to I,
 \qquad
 W_{R,S,-}^{[U]}\text{ being strong Cauchy},
\]
or the existence of a global Object~X is drawn here.
\end{openproblem}
